\lecture{2}{Sun 30 Jan 2022 15:50}{Relaxed SP}
\subsection{relaxed model PI}
\subsubsection{calculation 1 (in meeting)}

\begin{align*}
&\sum_{i=1}^n \prob{\text{accept }i \cap i \text{ best}}\\
&= \sum_{i=1}^n \prob{\text{accept }i | i \text{ best}}\prob{i \text{ best}}\\
&= \frac{1}{n}\sum_{i=1}^n \prob{X_1,\dots,X_{i-1}<T \cap X_i \geq T | i \text{ best}}\\
&= \frac{1}{n}\sum_{i=1}^n \prob{X_1,\dots,X_{i-1}<T | X_i \geq T \cap i \text{ best}}\prob{X_i \geq T | i\text{ is best}}\\ 
\end{align*}

Note that $\prob{X_i \geq T | i \text{ best}} = 1-F_{X_{\max}}(T) = 1-T^n$. Also,

\begin{align*}
&\prob{X_1,\dots,X_{i-1}<T | X_i \geq T \cap i \text{ best}}\\
&= \int_{T}^1 \prob{X_1,\dots,X_{i-1}<T | i \text{ best} \cap X_i = x}f_{X}(x | X_i \geq T \cap i \text{ best})\,dx\\
&= \int_{T}^1 \prob{X_1,\dots,X_{i-1}<T | X_1,\dots,X_{i-1}<x}f_{X_{\max}}(x | X_{\max}\geq T)\,dx\\
&= \int_{T}^1 {\left(\frac{T}{x}\right)}^{i-1} \frac{nx^{n-1}}{1-T^{n}} \,dx\\
&= \frac{nT^{i-1}}{1-T^{n}}\int_{T}^1  x^{n-i} \,dx\\
&= \frac{nT^{i-1}}{1-T^{n}} \frac{1-T^{n-i+1}}{n-i+1}
\end{align*}

Combining all terms,

\begin{align*}
&\sum_{i=1}^n \prob{\text{accept }i \cap i \text{ best}}\\
&=\frac{1}{n}\sum_{i=1}^n \frac{nT^{i-1}}{1-T^{n}} \frac{x^{n-i+1}}{n-i+1} \left(1-T^n\right)\\
&=\sum_{i=1}^n T^{i-1} \frac{1-T^{n-i+1}}{n-i+1}\\
\end{align*}
\subsubsection{Calculation 2}

We have $n$ items with weights $X_1, \dots, X_n \sim \text{Unif}(0,1)$.
Observed values are $\tilde X_i \sim \text{Unif}(X_i  - \delta, X_i + \delta)$.
We adopt a thresholding scheme for $\ALG$: We examine an item if $\tilde x_i > T$ and accept if $x_i > T$.

If the player reaches $X_n$, then $X_n$ must be accepted.

\begin{align*}
\begin{split}
\prob{X_{max} \text{ selected}} &=\sum_{i=1}^{ n -1}\prob{X_1, \dots, X_{i-1}<T \cap X_i \geq T \cap X_i \text{ is max}} \\
				&\qquad + \prob{X_1, \ldots, X_{n-1} < T \cap X_n \max}
\end{split} 
\\
\begin{split}
&=\sum_{i=1}^{ n -1}\int_0^1 \prob{X_1,\dots,X_{i-1}<T\cap X_i\geq T \cap X_{i+1},\dots,X_{n}\leq X_i | X_i = x} f_{\text{Unif}(0,1)}(x) dx \\
&\qquad +\frac{1}{n} \int_0^1 \prob{X_n<T | X_n<x} f_{\max \text{Unif(0,1)}}(x) \,dx
\end{split}
\\
\begin{split}
&=\sum_{i=1}^{ n -1}\int_T^1 \prob{X_1,\dots,X_{i-1}<T \cap X_{i+1},\dots,X_n\leq x} f_{\text{Unif}(0,1)}(x) dx \\
&\qquad + \ldots
\end{split}
\\
&=\sum_{i=1}^{ n -1}\int_T^1 T^{i-1} x^{n-i} dx + \int_0^T x^{n-1}\,dx + \int_T^1 x^{n-1} \frac{T}{x}\,dx
\\
&=\sum_{i=1}^{ n -1}\frac{1}{n-i+1} T^{i-1}(1-T^{n-i+1}) + \frac{1}{n}T^n + T\frac{ 1-T^{n-1} }{n-1}
\end{align*}
\section{Relaxed SP Prophet Inequality}
\subsection{Lower Bound of $P^{\text{max}}$}

In this section, we provide a lower bound of $P^{\text{max}}$, where
\[
P^{\text{max}} = \max_{T \in [0,1]}\sum_{i=1}^n\frac{1}{n-i+1} T^{i-1}(1-T^{n-i+1}).
\]

The optimal $T^{*}$ can be solved by the optimization problem. Let's show the lower bound of $P^{\text{max}}(T^{*})$. Let $T=\frac{n}{n+1}$, then we have
\begin{align*}
    P^{\text{max}}(T^{*}) &\geq P^{\text{max}}(T) \\&= \sum_{i=1}^n\frac{1}{n-i+1} T^{i-1}(1-T^{n-i+1}) \\&= \sum_{i=1}^n\frac{1}{n-i+1} T^{i-1}- \sum_{i=1}^n\frac{1}{n-i+1} T^{n} \\&= T^n \sum_{i=1}^n\frac{1}{i} (T^{-i}-1) 
\end{align*}

We take the limit $n \to \infty$, we have
\begin{align*}
    \lim_{n \to \infty} T^n \sum_{i=1}^n\frac{1}{i} (T^{-i}-1)  &= \frac{1}{e} \lim_{n \to \infty} \sum_{i=1}^n\frac{1}{i} ((1+\frac{1}{n})^i-1)
\end{align*}

Observe that
\[
(1+\frac{1}{n})^i-1=\sum_{k=1}^{i}{i \choose k}\frac{1}{n^k}.
\]

Then, we take the sum on $i$, we have
\[
\sum_{i=1}^n \frac{1}{i}(1+\frac{1}{n})^i-1=\sum_{i=1}^n \frac{1}{i} \sum_{k=1}^{i}{i \choose k}\frac{1}{n^k}
\]

By Fubini Theorem, we have
\[
\sum_{i=1}^n \frac{1}{i} \sum_{k=1}^{i}{i \choose k}\frac{1}{n^k} = \sum_{k=1}^n \sum_{i=k}^{n}\frac{1}{i}{i \choose k}\frac{1}{n^k}
\]

We have
\[
\frac{1}{i}{i \choose k} = \frac{(i-1)(i-2)...(i-k+1)}{k!}
\]

The nominator is a polynomial with order $k-1$ with respect to $i$. 

\begin{exercise}
\label{ex1}	
{\color{red}  Show that $\sum_{i=1}^ni^k = \frac{1}{k+1}n^{k+1}+O(n^k)$}
\end{exercise}

By this property, we have 
\[
\lim_{n \to \infty} \sum_{i=k}^{n}\frac{1}{i}{i \choose k}\frac{1}{n^k} = \frac{1}{k \cdot k!}\frac{n^k}{n^k} = \frac{1}{k \cdot k!}.
\]

Finally, we have the lower bound of $P^{\text{max}}$ is 
\[
\frac{1}{e}\sum_{k=1}^{n}\frac{1}{k \cdot k!} \approx 0.485.
\]

\subsection{Upper Bound of $P^{\text{max}}$}

Consider there are only two secretaries in the arriving sequence. We have $X_1 \sim \text{Unif}(0,1)$, $X_2 \sim \text{Unif}(0,1)$. 

\begin{exercise}
\label{ex2}	
{\color{red} Show that the upper bound of $P^{\text{max}}$ among all deterministic algorithms is $\frac{1}{2}$}
\end{exercise}

{\color{blue} Hint: Any deterministic algorithm can only choose $X_1$ or $X_2$. If $X_1=x$, what will happen? Then take integral. Another method is that think about a coupling random arrival instance $Y_1=X_1$, $Y_2=\text{Unif}(0,1)$. $Y$ obviously follows the arrival process distribution. When the first customer arrives, any algorithm cannot distinguish the arrival instance is $X$ or $Y$. Then for the first customer, if the algorithm accepts him, ...., if not accept,....}


\begin{exercise}
\label{ex3}
{\color{red}  Show that the upper bound of $P^{\text{max}}$ among all randomized algorithms is $\frac{1}{2}$}
\end{exercise}

{\color{blue} Hint: Yao's principle. (or set a probability distribution on $X_1$ and $X_2$)}.













